\section{Theorem B: single-sweep terminal landings}\label{sec:terminal-landings}

The class of this section has two designated families: an older family $e$, whose
lower region carries the only raw sweep, and a newer family $g$. We first count orbit
points of a $\mathrm Q_2$ selected core, then bound the frozen point area at one
selection and phase, and finally combine two selections with Theorem~G.

\subsection{Successful core words}

\begin{lemma}[Successful core words]\label{lem:core-words}
\formal{GapAreaMemo/CoreWords.lean}{SourceCovers.coreSuccess_prefix_or,SourceCovers.coreSuccess_internal_nonleaf,SourceCovers.core_root_count}
Let $\mathcal K$ be a $\mathrm Q_2$ set of complete families, and consider words of
families of $\mathcal K$.
\begin{enumerate}[label=(\alph*)]
\item Two words that are successful from the same time and have the same first letter
are comparable: one is a prefix of the other.
\item Every letter of a successful word other than its first and its last letter is
an original nonleaf of the overlap forest of $\mathcal K$.
\item Let $r$ be a time, $X$ a finite set of times and $\Phi\subseteq\mathcal K$.
Suppose that each $v\in X$ is the endpoint of a nonempty word $w_v$ that is
successful from $r$, and that the last letter of $w_v$ belongs to $\Phi$ whenever it is
a nonleaf. Then $|X|\le2|\Phi|+4$.
\end{enumerate}
Part (b) holds for every set of complete families.
\end{lemma}

\begin{outline}
(a) Induct on the words. After a nonidentity step of a family, the new time lies in
its span at or before its owner. That family no longer acts there, no other family of
its class acts there, and at most one family of the other class does
\cite[proof of Lemma~2.6]{SourceCovers2026}; the formal statement is
\lean{SourceCovers.core_unique_after}. So the second letters agree whenever both words
continue.

(b) Consecutive letters of a successful word have intersecting spans and strictly
decreasing owners, so the word is a simple path of the overlap forest
\cite[Lemma~2.6]{SourceCovers2026}. An internal letter has two distinct neighbours on
that path.

(c) Consider first the points whose words end in a nonleaf. Two such words with the
same first and last letters are comparable by (a), and the letters of a successful word
are distinct, so the words are equal. The first letter acts at $r$, and at most two
families act there; this gives at most $2|\Phi|$ points. For words that end in a leaf,
two words of length at least two with the same first letter are comparable by (a), and
a proper prefix would make its last letter, a leaf, an internal letter of the longer
word, contrary to (b). Recording the first letter and whether the length is one gives
at most four such points.
\end{outline}

\subsection{A newer family and a terminal older family}

\begin{lemma}[Adding a newer and a terminal older family]\label{lem:landing-orbits}
\formal{GapAreaMemo/LandingOrbits.lean}{SourceCovers.shadowOrbit_insert_newer,SourceCovers.shadowOrbit_insert_terminal,SourceCovers.raw_landing_antitone}
Let $\mathcal F$ be a set of complete families, $g$ and $e$ complete families, and $s$ a
time.
\begin{enumerate}[label=(\alph*)]
\item Every point of $\mathcal O_{\mathcal F\cup\{g\}}(s)$ lies in
$\mathcal O_{\mathcal F}(s)$, or in $\mathcal O_{\mathcal F}(g(u))$ for some
$u\in\mathcal O_{\mathcal F}(s)$ at which $g$ acts.
\item Suppose that no family of $\mathcal F$ acts at any time of the lower region
$(p_e,o_e]$. Then every point of $\mathcal O_{\mathcal F\cup\{e\}}(s)$ lies in
$\mathcal O_{\mathcal F}(s)$ or equals $e(u)$ for some $u\in\mathcal O_{\mathcal F}(s)$
at which $e$ acts.
\item If $x$ is nondecreasing on $(p_e,o_e]$, then $x(e(u'))\le x(e(u))$ whenever
$o_e<u\le u'\le U_e$.
\end{enumerate}
\end{lemma}

\begin{outline}
Parts (a) and (b) follow by induction over the letters of a word. (a) After a
nonidentity step of $g$ the time is at most $o_g$, and no later step increases it, so
$g$ never acts again. (b) After a nonidentity step of $e$ the time lies in
$(p_e,o_e]$ by \eqref{eq:strict-shadow}, where neither $e$ nor a family of
$\mathcal F$ acts; the point is final. (c) On its active interval the shadow of $e$ is
antitone with values in $(p_e,o_e]$; compose with the nondecreasing raw path.
\end{outline}

\subsection{Landings at one selection and phase}

\begin{definition}[Newer closure, landings and starts]\label{def:landings}
\formal{GapAreaMemo/LandingCounts.lean}{SourceCovers.newerRoots,SourceCovers.newerClosure,SourceCovers.landingSet,SourceCovers.IsLandingStep,SourceCovers.stepJump,SourceCovers.stepMax,SourceCovers.landingStarts}
Fix a $\mathrm Q_2$ selected core $\mathcal K$, $M\ge2$, a residue $\rho<M$, complete
families $g$ and $e$, and a time $s$. The \emph{newer roots} of $s$ are
$N_\rho(s)=\{s\}\cup\{g(u):u\in A_\rho(s),\ g(u)\ne u\}$, and its \emph{newer closure}
is $C_\rho(s)=\bigcup_{r\in N_\rho(s)}A_\rho(r)$. Since orbits are closed under
composition, this is the set $C_g(s)$ of \cite[Proposition~3.7(b)]{SourceCovers2026},
with $g$ in the role of the newer exceptional map. Its \emph{landing times} are
$\Lambda_\rho(s)=C_\rho(s)\cap(o_e,U_e]$. A \emph{landing step} is a pair $(u,F)$ of a
landing time $u$ and a surviving family $F$ that acts at $u$ with
$F(u)\in(o_e,U_e]$. Its \emph{jump} is $x(e(F(u)))-x(e(u))$, truncated at zero, and
$\Gamma_\rho(s)$ is the largest jump of a landing step, or zero if there is none.
The \emph{starts} $\Lambda^0_\rho(s)$ are the landing times $v$ for which there is no
landing time $u$ and surviving family $F$ acting at $u$ with $F(u)=v$.
\end{definition}

\begin{lemma}[Counts at one selection and phase]\label{lem:landing-counts}
\formal{GapAreaMemo/LandingCounts.lean}{SourceCovers.newerClosure_count,SourceCovers.newerClosure_count_span,SourceCovers.landingStarts_card,SourceCovers.landing_cover}
In the setting of \cref{def:landings}:
\begin{enumerate}[label=(\alph*)]
\item If $D$ is a set of times and $\Phi\subseteq\mathcal K$ contains every original
nonleaf whose lower region meets $D$, then $|C_\rho(s)\cap D|\le(2M+4)(2|\Phi|+5)$.
\item If $F\in\mathcal K$, $p_F<a$ and $v\le U_F$, then
$|C_\rho(s)\cap(a,v)|\le11(2M+4)$.
\item $|\Lambda^0_\rho(s)|\le5(2M+4)$.
\item If $x$ is nondecreasing on $(p_e,o_e]$, then every $v\in\Lambda_\rho(s)$ has a
start $s_0\in\Lambda^0_\rho(s)$ with
\[
 x(e(s_0))\le x(e(v))\le x(e(s_0))+(M+1)\,\Gamma_\rho(s).
\]
\end{enumerate}
\end{lemma}

\begin{outline}
(a) There are at most $1+|A_\rho(s)|\le2M+4$ newer roots. Fix one root $r$. Erasing
inactive letters \cite[Lemma~2.6]{SourceCovers2026}, every point of $A_\rho(r)$ other
than $r$ is the endpoint of a nonempty successful word of surviving families from $r$.
By \eqref{eq:strict-shadow} the endpoint lies in the lower region of the last letter,
so for an endpoint in $D$ a nonleaf last letter belongs to $\Phi$.
\Cref{lem:core-words}(c) leaves at most $2|\Phi|+4$ such points, and with $r$ at most
$2|\Phi|+5$ points per root.

(b) Take for $\Phi$ the family $F$ together with the nonleaves adjacent to $F$. A
nonleaf whose lower region meets $(a,v)\subseteq(p_F,U_F]$ is $F$ or has a span
meeting that of $F$. A nonleaf has at most two nonleaf neighbours and a leaf has at
most one neighbour, so $|\Phi|\le3$, and (a) gives $11(2M+4)$.

(c) Every start is a newer root, or lies in $A_\rho(r)$ for a newer root $r$ and
differs from it. In the second case it is the endpoint of a nonempty successful
survivor word from $r$. Its last step starts at a point of the newer closure that is
larger than the start and is not a landing time, since otherwise the start would be
reached by a landing step; so that point exceeds $U_e$. The span of the last letter
therefore contains the boundary between $U_e$ and $U_e+1$, which at most two families
of $\mathcal K$ do. The first letter acts at $r$, so it is one of at most two
families. Two such words with the same first and last letters are equal, by
\cref{lem:core-words}(a) and the distinctness of the letters of a successful word.
Each root thus contributes at most $1+4$ starts.

(d) Use strong induction on $U_e-v$. A landing time $v$ that is not a start equals
$F(u)$ for a landing time $u>v$ and a surviving family $F$; then $(u,F)$ is a landing
step, which raises the raw value by at most $\Gamma_\rho(s)$. The chain of steps from
a start $s_0$ to $v$ is a successful word of surviving families, so it has at most
$M+1$ letters \cite[Lemma~2.6]{SourceCovers2026}. The lower bound is
\cref{lem:landing-orbits}(c), since $v\le s_0$.
\end{outline}

\subsection{Single-sweep terminal-landing geometries}

\begin{definition}[Single-sweep terminal-landing geometry]\label{def:terminal-landing}
\formal{GapAreaMemo/LandingGaps.lean}{SourceCovers.AdmissibleGeometry.TerminalLanding}
Single-sweep terminal-landing data for an admissible geometry, with original core
$\mathcal K_0$ and exceptional set $\mathcal E$, consist of complete families $e$ and
$g$, the \emph{older} and the \emph{newer} family, and a natural number $x_0$, such that
\begin{enumerate}[label=(L\arabic*)]
\item $\mathcal E\subseteq\{e,g\}$;
\item $x(t)=x_0$ whenever $t$ lies in the lower region of a family of $\mathcal K_0$ or
of $g$;
\item $x(s_j)=x_0$ for every column $j<d$;
\item $x$ is nondecreasing on the lower region $(p_e,o_e]$ of $e$;
\item no family of $\mathcal K_0$, and not $g$, acts at any time of $(p_e,o_e]$.
\end{enumerate}
A geometry with such data is a \emph{single-sweep terminal-landing geometry}.
The families $e$ and $g$ need not be distinct, and need not both be exceptional.
Condition (L5) holds for $g$ whenever $o_e\le o_g$ (an elementary remark, not
formalized).
\end{definition}

Thus the input is read in one nondecreasing sweep during the lower region of $e$,
the input head is parked at $x_0$ at all other times that can contribute to a frozen
bank (\cref{lem:landing-gaps}(a)), and no family of a bank moves a landing of $e$
again.

\begin{lemma}[Finite gap bounds for terminal landings]\label{lem:landing-gaps}
\formal{GapAreaMemo/LandingGaps.lean}{SourceCovers.AdmissibleGeometry.TerminalLanding.raw_values,SourceCovers.AdmissibleGeometry.TerminalLanding.gap_of_step,SourceCovers.AdmissibleGeometry.TerminalLanding.gap_of_starts}
Let an admissible geometry carry single-sweep terminal-landing data $(e,g,x_0)$. Let
$\mathcal K\subseteq\mathcal K_0$ be a $\mathrm Q_2$ selected core, $M\ge2$, $\rho<M$ a
residue and $j<d$ a column. Landing sets refer to $\mathcal K$, $M$, $\rho$, $g$, $e$
and the root $s_j$, and $S_\rho(s_j)$ to the geometry's exceptional set.
\begin{enumerate}[label=(\alph*)]
\item Every element of $S_\rho(s_j)$ is $x_0$ or $x(e(v))$ for some
$v\in\Lambda_\rho(s_j)$.
\item Let $F\in\mathcal K$ act at a time $u<T$ with $o_e<u\le U_e$ and
$o_e<F(u)\le U_e$, and let $\gamma=x(e(F(u)))-x(e(u))$, truncated at zero. Then
\[
 \gaparea(S_\rho(s_j))\ge\frac{(\gamma/L)^2}{2\bigl(11(2M+4)+3\bigr)}.
\]
The family $F$ need not survive at $\rho$.
\item With $\Gamma=\Gamma_\rho(s_j)$,
\[
 \gaparea(S_\rho(s_j))\ge
 \frac{\bigl(1-5(2M+4)(M+1)\Gamma/L\bigr)_+^2}{2\bigl(5(2M+4)+1\bigr)}.
\]
\end{enumerate}
\end{lemma}

\begin{outline}
(a) By (L1), the orbit $O_\rho(s_j)$ of the retained core with the exceptions lies in
the orbit of the retained core together with $g$ and $e$. By (L5) and
\cref{lem:landing-orbits}(b), each of its points lies in the orbit of the retained
core with $g$, or equals $e(u)$ for a point $u$ of that orbit at which $e$ acts. A
point of the first kind is the root, where $x=x_0$ by (L3), or lies in the lower
region of a retained family or of $g$, where $x=x_0$ by (L2). For the second kind,
$u\in C_\rho(s_j)$ by \cref{lem:landing-orbits}(a), and $u\in(o_e,U_e]$ because $e$
acts at $u$.

(b) Since $F(u)<u$, \cref{lem:landing-orbits}(c) gives $x(e(u))\le x(e(F(u)))$. By (a)
and the same monotonicity, every raw value strictly between them is $x_0$ or
$x(e(v))$ for a landing time $v$ with $F(u)<v<u$. Since $F$ acts at $u$, $p_F<F(u)$
and $u\le U_F$, so \cref{lem:landing-counts}(b) leaves at most $1+11(2M+4)$ raw values
in that open interval. Cover the raw set by the two complementary closed intervals and
these singletons, and apply \eqref{eq:interval-budget}; the formal step is
\lean{SourceCovers.pointArea_ge_of_sparse}.

(c) By (a) and \cref{lem:landing-counts}(d), the raw set lies in $\{x_0\}$ and the
closed intervals $[x(e(s_0)),x(e(s_0))+(M+1)\Gamma]$ with $s_0\in\Lambda^0_\rho(s_j)$.
By \cref{lem:landing-counts}(c) these are at most $1+5(2M+4)$ intervals of total length
at most $5(2M+4)(M+1)\Gamma$. Apply \eqref{eq:interval-budget}.
\end{outline}

\subsection{Theorem B}

\begin{lemma}[Prefix bounds for terminal landings]\label{lem:landing-prefix}
\formal{GapAreaMemo/TerminalLandings.lean}{SourceCovers.AdmissibleGeometry.landingGamma,SourceCovers.AdmissibleGeometry.landingEnergy,SourceCovers.AdmissibleGeometry.landingEnergy_bounds,SourceCovers.AdmissibleGeometry.prefixMeanGapArea_ge_later,SourceCovers.AdmissibleGeometry.prefixMeanGapArea_ge_own}
Let an admissible geometry carry single-sweep terminal-landing data, and fix an order
of its original core by nonincreasing hop count. For $J_0\ge0$ and a column $j$, let $\Gamma_{J_0}(j)$ be the largest of the
jumps $\Gamma_\rho(s_j)$ over all selections of the first $J\le J_0$ original core
families with the first $R\le J$ protected and modulus $2R+2$, and over the legal
phases $\rho$ of each selection. The \emph{landing energy}
\[
 E(J_0)=\frac1d\sum_{j<d}\Bigl(\frac{\Gamma_{J_0}(j)}L\Bigr)^2
\]
lies in $[0,1]$. Put $f_R=20(R+2)(2R+3)$. For every $R\ge0$:
\begin{enumerate}[label=(\alph*)]
\item $\meanarea_{J,R,2R+2}\ge E(J_0)/(88R+182)$ for every $J\ge J_0$;
\item $\meanarea_{J,R,2R+2}\ge\bigl(1/2-2f_R^2E(J)\bigr)/\bigl(2(20R+41)\bigr)$ for every
$J\ge R$.
\end{enumerate}
\end{lemma}

\begin{outline}
At $M=2R+2$ the constants are $88R+182=2(11(2M+4)+3)$, $20R+41=5(2M+4)+1$ and
$f_R=5(2M+4)(M+1)$. Each jump is at most $L$, so $E(J_0)\le1$.

(a) Either $\Gamma_{J_0}(j)=0$, or it is the jump of an actual landing step $(u,F)$ at
a selection of the first $J'\le J_0$ families. Then $F$ also belongs to the first
$J\ge J_0$ families, and \cref{lem:landing-gaps}(b) at the selection $(J,R,2R+2)$ and
any of its legal phases gives
$\gaparea(S_\rho(s_j))\ge(\Gamma_{J_0}(j)/L)^2/(88R+182)$; the family need not survive
there. Average over legal phases and columns.

(b) For $J\ge R$, the selection $(J,R,2R+2)$ is among those that define
$\Gamma_J(j)$, so $\Gamma_\rho(s_j)\le\Gamma_J(j)$ at each of its legal phases. With
$y=\Gamma_J(j)/L$ and $\alpha=f_R\Gamma_\rho(s_j)/L\le f_Ry$,
\[
 (1-\alpha)_+^2\ge1-2\alpha\ge1-2f_Ry\ge\tfrac12-2f_R^2y^2,
\]
the last step because $2(f_Ry-\tfrac12)^2\ge0$. Apply \cref{lem:landing-gaps}(c) and
average over phases and columns.
\end{outline}

\begin{theorem}[Theorem B: single-sweep terminal landings are excluded]\label{thm:terminal-landings}
\formal{GapAreaMemo/TerminalLandings.lean}{SourceCovers.NearFullSequence.not_eventually_terminalLanding,SourceCovers.NearFullSequence.not_frequently_terminalLanding}
In no near-full sequence are all geometries from some index on single-sweep
terminal-landing geometries, and in none are infinitely many of its geometries such
geometries. The parameters $c,\beta,N,q$ of the sequence are arbitrary, and no envelope
on the sizes of the core families is assumed.
\end{theorem}

\begin{outline}
Fix a joint profile; along its subsequence the geometries eventually carry such data.

\emph{Step 1: the landing energies vanish.} For every $J_0$, the landing energy of the
geometry of index $\nu$ tends to zero along the subsequence. Otherwise it is at least
some $\eps>0$ infinitely often. By \cref{lem:landing-prefix}(a), for every $R$ and
every $J\ge\max(J_0,R)$ the inner upper limit of $\meanarea_{\nu,J,R,2R+2}$ is at least
$\eps/(88R+182)$, so $\limarea_R\ge\eps/(88R+182)$. Theorem~G gives
$\eps/(88R+182)\le\beta\tau_R/(R+2)$, and since $88R+182\le91(R+2)$, $\eps\le91\beta\tau_R$
for every $R$, contrary to $\tau_R\to0$.

\emph{Step 2: the selections themselves.} Fix $R$ and $J\ge R$. By Step~1, eventually
$2f_R^2E(J)\le1/4$, and \cref{lem:landing-prefix}(b) gives
$\meanarea_{\nu,J,R,2R+2}\ge1/(8(20R+41))$. Hence $\limarea_R\ge1/(8(20R+41))$, and
Theorem~G gives
\[
 \beta\tau_R\ge\frac{R+2}{8(20R+41)}\ge\frac1{164}\qquad(R\ge0),
\]
contrary to $\tau_R\to0$. If such geometries occur infinitely often, they form a
near-full subsequence, and the first case applies to it.
\end{outline}

\paragraph{Why it matters.}
The two-exception class with an inverse-square hop envelope is the nearest unresolved
case of \cite[Section~9.1]{SourceCovers2026}, and the necessary hop-tail bound
\cite[Theorem~4.2]{SourceCovers2026} is compatible with it. Theorem~B excludes, with no
envelope, the two-exception geometries in which the older family carries the only raw
sweep and receives terminal landings. A near-full sequence whose geometries have at
most two exceptional families must therefore, from some index on and for every choice
of $e$, $g$ and $x_0$ with $\mathcal E\subseteq\{e,g\}$, violate one of (L2)--(L5): it reads a coordinate
other than $x_0$ during the lower region of a core family or of $g$, or at a column
root; or its raw path decreases somewhere on the lower region of $e$; or a core family
or $g$ acts during that lower region. For an impossibility proof the argument shows
how Theorem~G is used at two scales: a fixed prefix bounds the landing jumps of every
larger selection (\cref{lem:landing-prefix}(a)), and at the selection itself few
starts with small jumps confine the frozen raw set to a few short intervals
(\cref{lem:landing-prefix}(b)).
